\documentclass[conference]{IEEEtran}
\IEEEoverridecommandlockouts
% The preceding line is only needed to identify funding in the first footnote. If that is unneeded, please comment it out.
\usepackage{cite}
\usepackage{amsmath,amssymb,amsfonts}
\usepackage{algorithmic}
\usepackage{graphicx}
\usepackage{textcomp}
\usepackage{xcolor}
\usepackage{booktabs}
\usepackage{url}
\usepackage{tikz}
\usetikzlibrary{positioning,arrows.meta,shapes.geometric,fit,backgrounds}
\def\BibTeX{{\rm B\kern-.05em{\sc i\kern-.025em b}\kern-.08em
    T\kern-.1667em\lower.7ex\hbox{E}\kern-.125emX}}
\begin{document}

\title{Order-2 Co-occurrence TF-IDF for Interaction-Free Fragrance Dupe Retrieval: A Non-Circular Study}

\author{\IEEEauthorblockN{Azam Azri}
\IEEEauthorblockA{\textit{Master's Program in Informatics} \\
\textit{Affiliation (to be completed)}\\
City, Indonesia \\
azamazrri@gmail.com}
}

\maketitle

\begin{abstract}
Content-based perfume recommenders predominantly rank candidates by order-one
accord similarity, such as cosine or Jaccard over note sets. Whether modeling the
co-occurrence of accords improves retrieval has not been tested under a leakage-free
protocol, and the prevalent catalogs carry no user interactions, which rules out
collaborative and knowledge-graph recommenders that require an interaction signal.
This work studies interaction-free fragrance dupe retrieval: given a global perfume,
retrieve the local product that imitates it, with the held-out inspired-by edge as
ground truth. A parameter-free matcher is proposed that augments accord TF-IDF with
accord-pair (order-two) tokens; its cosine score decomposes additively into
order-one and order-two terms and yields faithful, auditable co-occurrence paths for
explanation. Evaluation uses 340 products and 209 queries under leave-one-out
ranking, with a three-source firewall whose query accords are independently verified
from Fragrantica. Against thirteen methods spanning marginal similarity, neural
embeddings, perceptual-taxonomy transport, and learned models, the proposed method
attains the best mean reciprocal rank (MRR 0.496) and beats the TF-IDF baseline
(0.454) by a modest but robust margin (delta-MRR +0.042, 95\% CI [+0.019, +0.067],
Wilcoxon p = 8.7e-8); order-two terms contribute 81.3\% of the top-ranked score.
The effect is honestly modest, and limitations are discussed.
\end{abstract}

\begin{IEEEkeywords}
recommender systems, content-based retrieval, co-occurrence, TF-IDF, explainability,
fragrance, non-circular evaluation
\end{IEEEkeywords}

\section{Introduction}
Selecting a fragrance ``dupe'' --- an affordable local product engineered to imitate
the scent of a well-known global perfume --- is a common consumer task and a natural
content-based retrieval problem. Each product is described by a small set of
olfactory \emph{accords} (e.g., \textit{amber}, \textit{woody}, \textit{warm spicy}),
and a dupe is expected to share much of the accord profile of the global perfume it
imitates. Existing perfume recommenders rank candidates almost exclusively by
order-one accord similarity, typically cosine over TF-IDF vectors or Jaccard over
accord sets \cite{r23,r24,r28}. Such marginal representations treat accords
independently and discard the structure in which accords co-occur.

A second constraint is decisive in practice: many real product catalogs record only
content (accord lists and family labels) and carry \emph{no} user interactions or
ratings. This excludes by construction the interaction-driven recommenders that
dominate the higher-order and knowledge-graph literature --- collaborative filtering,
LightGCN, KGIN, and KGAT \cite{r20} --- because they require a user--item signal the
data does not contain. Methods that do exploit item co-occurrence and higher-order
structure report strong gains, but only in interaction-based settings such as
session or multi-behavior recommendation \cite{r1,r2,r3,r4,r5}. Whether order-two
accord co-occurrence helps in a purely content-based, interaction-free regime, and
whether it does so under an evaluation that cannot leak the label, remains open.

A third assumption pervades adjacent fields: that injecting perceptual or learned
\emph{structure} aids retrieval. Perceptual odor science organizes accords on sensory
wheels and perceptual maps \cite{r10,r12,r13}, and a large body of work assumes such
geometry should help. Yet evidence from small and imbalanced datasets repeatedly
shows that simple baselines match or beat flexible models that overfit
\cite{r14,r15,r16,r17}. This tension has not been examined for interaction-free
fragrance retrieval.

This paper addresses these gaps with a deliberately simple, parameter-free method and
a rigorous, non-circular protocol. Accords and accord-pairs (unigrams and bigrams of
a co-occurrence graph) are weighted by inverse document frequency and compared by
cosine; the score decomposes additively into an order-one (marginal) term and an
order-two (co-occurrence) term, and the shared accord-pairs form auditable
explanation paths grouped by fragrance-wheel super-family. The matcher is intentionally
free of trained parameters: added capacity is shown empirically to overfit. The
contribution is a \emph{finding and a protocol}, not a new model, calibrated to a
moderate, applied scope.

The main contributions are:
\begin{itemize}
\item \textbf{Task and protocol.} We formalize interaction-free dupe retrieval and a
non-circular evaluation with a three-source firewall (local item accords, externally
sourced query accords verified from Fragrantica, and held-out inspired-by labels), and
we make explicit why interaction-based recommenders (KGAT, LightGCN, KGIN) cannot be
compared on this data.
\item \textbf{Method.} An order-two co-occurrence TF-IDF matcher whose cosine score
decomposes into order-one and order-two terms, paired with faithful-by-construction
co-occurrence explanation paths for a downstream explainability layer.
\item \textbf{Finding.} A first non-circular demonstration that order-two
co-occurrence is the strongest of fourteen methods spanning four paradigms (marginal,
neural embedding, perceptual taxonomy, and learned), reported with honest caveats: the
gain over the TF-IDF baseline is modest (+0.042 MRR) but robust and significant; the
Sentence-BERT loss is partly attributable to query-product text asymmetry; and the
moderate learned methods only match, rather than fall far below, the baseline.
\item \textbf{Reproducibility.} A complete, deterministic benchmark of fourteen
methods with paired significance testing and bootstrap confidence intervals.
\end{itemize}

The remainder of the paper reviews related work; formalizes the problem and the
non-circular firewall; presents the proposed method; describes the experimental setup;
reports results; discusses them; and states limitations before concluding.

\bibliographystyle{IEEEtran}
\bibliography{references}

\end{document}
